\chapter{Patterns}\label{record-023}

\noindent{\footnotesize\color{BookMuted}\path{docs/design-system/patterns.md}\par}\smallskip

This document defines the recurring product patterns that Bijou should support consistently.

These are the decision frameworks behind the component families.

\section{0. Selection, focus, and rhythm}\label{record-023-selection-focus-and-rhythm}

\subsection{The pattern}\label{record-023-the-pattern}

Bijou should make it obvious:

\begin{itemize}
\tightlist
\item
  which item is currently selected
\item
  which region currently owns input
\item
  which copy is primary versus supporting
\item
  how dense content is meant to breathe
\end{itemize}

These decisions should not be improvised per surface.

\subsection{Canonical rules}\label{record-023-canonical-rules}

\subsubsection{Background fill means current selection}\label{record-023-background-fill-means-current-selection}

Use background fill to indicate the current selected row or item inside a collection.

Why:

\begin{itemize}
\tightlist
\item
  it reads immediately in dense TUIs
\item
  it does not depend on one tiny pointer glyph
\item
  it survives better than hue-only emphasis when terminals differ
\end{itemize}

Selection may also use supporting cues such as:

\begin{itemize}
\tightlist
\item
  a leading chevron
\item
  bold value text
\item
  a stronger border inside the selected row
\end{itemize}

But those are supporting cues. The canonical selected-item treatment is still background fill.

\subsubsection{Structural accent means active region}\label{record-023-structural-accent-means-active-region}

Use a structural accent such as a gutter, border emphasis, or region frame to show the active region.

Why:

\begin{itemize}
\tightlist
\item
  region focus and row selection are not the same thing
\item
  a pane can be active while its selected row changes
\item
  a selected row can exist inside an inactive region without stealing shell ownership
\end{itemize}

The accent should describe ownership of input, not merely visual importance.

\subsubsection{Visible controls are a promise}\label{record-023-visible-controls-are-a-promise}

Follow \href{https://github.com/flyingrobots/bijou/blob/cbd70612140e7a450d764b3c7be1c6af878b99bc/docs/invariants/visible-controls-are-a-promise.md}{Visible Controls Are a Promise}.

If a region is not active, its controls should not be advertised as though they will work.

When the shell owns focus, shell controls replace pane-local cues.

\subsubsection{One cell of padding matters}\label{record-023-one-cell-of-padding-matters}

Dense TUIs still need rhythm.

Canonical defaults:

\begin{itemize}
\tightlist
\item
  one cell of padding between content and region edges
\item
  one blank line between titled sections
\item
  one blank line between preference rows or similar stacked interactive rows when the surface can afford it
\end{itemize}

Do not cram labels, values, and descriptions together when one cell of space would make the surface easier to parse.

\subsubsection{Stack beneath the label before truncating}\label{record-023-stack-beneath-the-label-before-truncating}

When labels and values compete for the same horizontal space:

\begin{itemize}
\tightlist
\item
  wrap or stack first
\item
  reflow second
\item
  truncate only when the surface truly demands single-line identity
\end{itemize}

For settings and similar rows, values should stack beneath the label before truncating into ambiguity.

Marquee is an explicit overflow behavior for compact, focused single-line surfaces. It is not the default solution.

\subsubsection{Supporting copy should look supporting}\label{record-023-supporting-copy-should-look-supporting}

Descriptions, helper copy, and secondary explanation should look visibly quieter than primary actions and current values.

Use:

\begin{itemize}
\tightlist
\item
  muted tokens
\item
  dim modifiers when needed
\item
  more breathing room around primary actions than around support prose
\end{itemize}

Do not make supporting copy compete with what the user is supposed to act on now.

\subsubsection{Layers should behave like a stack}\label{record-023-layers-should-behave-like-a-stack}

When interactive layers overlap, they should behave in top-down order.

Canonical rules:

\begin{itemize}
\tightlist
\item
  the topmost layer owns input
\item
  the topmost layer owns the visible control hints
\item
  dismiss actions target the topmost dismissible layer first
\item
  underlying regions do not keep advertising controls while covered
\end{itemize}

This means a palette, drawer, or modal should not merely look on top. It should also be the source of truth for routing and footer/help copy.

\section{1. Status and feedback}\label{record-023-status-and-feedback}

\subsection{The pattern}\label{record-023-the-pattern-1}

Not all feedback should interrupt the user equally. Choose the least disruptive surface that still preserves the needed meaning.

\subsection{Escalation ladder}\label{record-023-escalation-ladder}

\subsubsection{\texorpdfstring{\texttt{b\allowbreak{}a\allowbreak{}d\allowbreak{}g\allowbreak{}e\allowbreak{}(\allowbreak{})}}{}}\label{record-023-badge}

Use when:

\begin{itemize}
\tightlist
\item
  the status is compact
\item
  it belongs inline with another object
\item
  the surrounding content already provides context
\end{itemize}

Avoid when:

\begin{itemize}
\tightlist
\item
  the user must act on the information
\item
  the message needs explanation
\item
  the status must persist as its own block
\end{itemize}

\subsubsection{\texorpdfstring{\texttt{n\allowbreak{}o\allowbreak{}t\allowbreak{}e\allowbreak{}(\allowbreak{})}}{}}\label{record-023-note}

Use when:

\begin{itemize}
\tightlist
\item
  the user needs supporting explanation
\item
  the content is informational, not urgent
\item
  the message belongs inside a form or guided flow
\end{itemize}

Avoid when:

\begin{itemize}
\tightlist
\item
  urgency or danger must be communicated
\item
  the message should interrupt
\end{itemize}

\subsubsection{\texorpdfstring{\texttt{a\allowbreak{}l\allowbreak{}e\allowbreak{}r\allowbreak{}t\allowbreak{}(\allowbreak{})}}{}}\label{record-023-alert}

Use when:

\begin{itemize}
\tightlist
\item
  the state should remain in the document flow
\item
  the message is important enough to deserve its own block
\item
  the user should see it while reading or editing nearby content
\end{itemize}

Avoid when:

\begin{itemize}
\tightlist
\item
  the event is brief and transient
\item
  it belongs in an activity stream
\item
  the app should manage stack/history/lifecycle
\end{itemize}

\subsubsection{\texorpdfstring{\texttt{t\allowbreak{}o\allowbreak{}a\allowbreak{}s\allowbreak{}t\allowbreak{}(\allowbreak{})}}{}}\label{record-023-toast}

Use when:

\begin{itemize}
\tightlist
\item
  the event is transient
\item
  the app is composing a one-off overlay directly
\item
  custom placement matters, but full notification lifecycle does not
\end{itemize}

Avoid when:

\begin{itemize}
\tightlist
\item
  the message must remain in the page content
\item
  the app needs stacking, history, routing, or recall
\item
  the user needs to review a set of related events later
\end{itemize}

Content guidance:

\begin{itemize}
\tightlist
\item
  keep the message short and self-contained
\item
  do not hide follow-up obligations in a toast body
\item
  prefer one clear state statement over mixed status plus instructions plus history
\end{itemize}

\subsubsection{notifications}\label{record-023-notifications}

Use when:

\begin{itemize}
\tightlist
\item
  the app owns transient messaging as a system
\item
  stacking, placement, actions, routing, or history matter
\item
  the user may need to reopen or review prior notices
\end{itemize}

Avoid when:

\begin{itemize}
\tightlist
\item
  the message should remain in page content
\item
  a single local overlay would do
\item
  the user must stop and decide before continuing
\end{itemize}

Content guidance:

\begin{itemize}
\tightlist
\item
  use notifications for events, not for durable documents
\item
  route warnings and errors here when the app owns runtime messaging as a system
\item
  actionable notifications should expose one obvious next step, not a mini workflow
\item
  if users may need recall, the history/archive surface is part of the feature, not an optional afterthought
\end{itemize}

\subsubsection{\texorpdfstring{\texttt{m\allowbreak{}o\allowbreak{}d\allowbreak{}a\allowbreak{}l\allowbreak{}(\allowbreak{})}}{}}\label{record-023-modal}

Use when:

\begin{itemize}
\tightlist
\item
  the user must stop and decide
\item
  the app needs an explicit review or confirmation surface
\item
  background actions should be blocked
\end{itemize}

Avoid when:

\begin{itemize}
\tightlist
\item
  a lighter inline or overlay treatment would do
\end{itemize}

Content guidance:

\begin{itemize}
\tightlist
\item
  the title should state the decision or review task clearly
\item
  the body should explain the consequence, not repeat generic warning copy
\item
  the modal should end in a clear next action or dismissal path
\item
  avoid turning a modal into a drawer with extra blocking chrome
\end{itemize}

\subsubsection{\texorpdfstring{\texttt{s\allowbreak{}k\allowbreak{}e\allowbreak{}l\allowbreak{}e\allowbreak{}t\allowbreak{}o\allowbreak{}n\allowbreak{}(\allowbreak{})}}{}}\label{record-023-skeleton}

Use when:

\begin{itemize}
\tightlist
\item
  the content shape is known but a short-lived loading gap would otherwise cause layout pop
\item
  the user benefits from seeing where content will land
\end{itemize}

Avoid when:

\begin{itemize}
\tightlist
\item
  honest partial content is already available
\item
  the loading state is long-lived enough that explicit progress or retry messaging would be more honest
\item
  the placeholder would imply certainty the app does not actually have yet
\end{itemize}

\subsubsection{\texorpdfstring{\texttt{p\allowbreak{}r\allowbreak{}o\allowbreak{}g\allowbreak{}r\allowbreak{}e\allowbreak{}s\allowbreak{}s\allowbreak{}B\allowbreak{}a\allowbreak{}r\allowbreak{}(\allowbreak{})} / \texttt{c\allowbreak{}r\allowbreak{}e\allowbreak{}a\allowbreak{}t\allowbreak{}e\allowbreak{}A\allowbreak{}n\allowbreak{}i\allowbreak{}m\allowbreak{}a\allowbreak{}t\allowbreak{}e\allowbreak{}d\allowbreak{}P\allowbreak{}r\allowbreak{}o\allowbreak{}g\allowbreak{}r\allowbreak{}e\allowbreak{}s\allowbreak{}s\allowbreak{}B\allowbreak{}a\allowbreak{}r\allowbreak{}(\allowbreak{})}}{ / }}\label{record-023-progressbar-createanimatedprogressbar}

Use when:

\begin{itemize}
\tightlist
\item
  percent-complete is actually known or can be estimated honestly
\item
  the user benefits from determinate progress instead of generic activity
\end{itemize}

Avoid when:

\begin{itemize}
\tightlist
\item
  completion is unknown and the percentage would be fake
\item
  the work is so brief that the bar would only flicker
\end{itemize}

Loading-screen rule:

\begin{itemize}
\tightlist
\item
  a loading screen should show at most one progress bar
\item
  that bar must represent actual loading progress
\item
  do not use a progress bar on a title or splash screen unless the screen is truly a loading surface
\end{itemize}

\subsubsection{\texorpdfstring{\texttt{s\allowbreak{}p\allowbreak{}i\allowbreak{}n\allowbreak{}n\allowbreak{}e\allowbreak{}r\allowbreak{}F\allowbreak{}r\allowbreak{}a\allowbreak{}m\allowbreak{}e\allowbreak{}(\allowbreak{})} / \texttt{c\allowbreak{}r\allowbreak{}e\allowbreak{}a\allowbreak{}t\allowbreak{}e\allowbreak{}S\allowbreak{}p\allowbreak{}i\allowbreak{}n\allowbreak{}n\allowbreak{}e\allowbreak{}r\allowbreak{}(\allowbreak{})}}{ / }}\label{record-023-spinnerframe-createspinner}

Use when:

\begin{itemize}
\tightlist
\item
  the task is active but indeterminate
\item
  the user needs reassurance that work is still happening
\end{itemize}

Avoid when:

\begin{itemize}
\tightlist
\item
  completion can be expressed honestly as a determinate bar
\item
  animation would distract more than it reassures
\end{itemize}

\section{2. Selection versus action}\label{record-023-selection-versus-action}

This is one of the easiest places for a component library to become fuzzy.

\subsection{\texorpdfstring{\texttt{s\allowbreak{}e\allowbreak{}l\allowbreak{}e\allowbreak{}c\allowbreak{}t\allowbreak{}(\allowbreak{})}}{}}\label{record-023-select}

Use when:

\begin{itemize}
\tightlist
\item
  the user is choosing one value
\item
  the set is short and stable
\item
  the result becomes stored state
\end{itemize}

Avoid when:

\begin{itemize}
\tightlist
\item
  the list is large enough that search is the real task
\item
  the outcome should trigger an immediate app command instead of persisting a value
\end{itemize}

\subsection{\texorpdfstring{\texttt{m\allowbreak{}u\allowbreak{}l\allowbreak{}t\allowbreak{}i\allowbreak{}s\allowbreak{}e\allowbreak{}l\allowbreak{}e\allowbreak{}c\allowbreak{}t\allowbreak{}(\allowbreak{})}}{}}\label{record-023-multiselect}

Use when:

\begin{itemize}
\tightlist
\item
  the user is building a set
\item
  multiple simultaneous selections are meaningful
\end{itemize}

Avoid when:

\begin{itemize}
\tightlist
\item
  the user should choose exactly one option
\item
  the user is really firing commands one at a time instead of building a lasting set
\end{itemize}

\subsection{\texorpdfstring{\texttt{f\allowbreak{}i\allowbreak{}l\allowbreak{}t\allowbreak{}e\allowbreak{}r\allowbreak{}(\allowbreak{})}}{}}\label{record-023-filter}

Use when:

\begin{itemize}
\tightlist
\item
  the user is choosing one value from a larger set
\item
  search/narrowing is the main task
\end{itemize}

Avoid when:

\begin{itemize}
\tightlist
\item
  the option set is already short and stable enough for \texttt{s\allowbreak{}e\allowbreak{}l\allowbreak{}e\allowbreak{}c\allowbreak{}t\allowbreak{}(\allowbreak{})}
\item
  the result is a command or navigation action rather than stored value state
\end{itemize}

\subsection{\texorpdfstring{\texttt{c\allowbreak{}o\allowbreak{}m\allowbreak{}m\allowbreak{}a\allowbreak{}n\allowbreak{}d\allowbreak{}P\allowbreak{}a\allowbreak{}l\allowbreak{}e\allowbreak{}t\allowbreak{}t\allowbreak{}e\allowbreak{}(\allowbreak{})} / \texttt{c\allowbreak{}o\allowbreak{}m\allowbreak{}m\allowbreak{}a\allowbreak{}n\allowbreak{}d\allowbreak{}P\allowbreak{}a\allowbreak{}l\allowbreak{}e\allowbreak{}t\allowbreak{}t\allowbreak{}e\allowbreak{}S\allowbreak{}u\allowbreak{}r\allowbreak{}f\allowbreak{}a\allowbreak{}c\allowbreak{}e\allowbreak{}(\allowbreak{})}}{ / }}\label{record-023-commandpalette-commandpalettesurface}

Use when:

\begin{itemize}
\tightlist
\item
  the result is an action, navigation, or command
\item
  the app wants one global action-discovery surface
\item
  the user may benefit from mixed command groups such as actions, navigation, or recent destinations
\end{itemize}

Content guidance:

\begin{itemize}
\tightlist
\item
  palette entries should read like things the app can do, not field values the app stores
\item
  group commands by user intent (\texttt{N\allowbreak{}a\allowbreak{}v\allowbreak{}i\allowbreak{}g\allowbreak{}a\allowbreak{}t\allowbreak{}i\allowbreak{}o\allowbreak{}n}, \texttt{A\allowbreak{}c\allowbreak{}t\allowbreak{}i\allowbreak{}o\allowbreak{}n\allowbreak{}s}, \texttt{R\allowbreak{}e\allowbreak{}c\allowbreak{}e\allowbreak{}n\allowbreak{}t}) instead of raw subsystem names when possible
\item
  keep descriptions short enough to scan without turning the palette into a browseable document
\end{itemize}

Avoid \texttt{c\allowbreak{}o\allowbreak{}m\allowbreak{}m\allowbreak{}a\allowbreak{}n\allowbreak{}d\allowbreak{}P\allowbreak{}a\allowbreak{}l\allowbreak{}e\allowbreak{}t\allowbreak{}t\allowbreak{}e\allowbreak{}(\allowbreak{})} / \texttt{c\allowbreak{}o\allowbreak{}m\allowbreak{}m\allowbreak{}a\allowbreak{}n\allowbreak{}d\allowbreak{}P\allowbreak{}a\allowbreak{}l\allowbreak{}e\allowbreak{}t\allowbreak{}t\allowbreak{}e\allowbreak{}S\allowbreak{}u\allowbreak{}r\allowbreak{}f\allowbreak{}a\allowbreak{}c\allowbreak{}e\allowbreak{}(\allowbreak{})} when:

\begin{itemize}
\tightlist
\item
  the user is filling out a field value
\item
  the result should behave like stored form data
\item
  the content is really a browseable record list rather than an action list
\end{itemize}

\section{3. Comparison versus browsing versus inspection}\label{record-023-comparison-versus-browsing-versus-inspection}

\subsection{\texorpdfstring{\texttt{t\allowbreak{}a\allowbreak{}b\allowbreak{}l\allowbreak{}e\allowbreak{}(\allowbreak{})}}{}}\label{record-023-table}

Use when:

\begin{itemize}
\tightlist
\item
  the user needs comparison across columns
\item
  the same attributes matter across many rows
\item
  the table should fit the available terminal width and wrap long cells instead of forcing callers to pre-compute column widths
\end{itemize}

Avoid when:

\begin{itemize}
\tightlist
\item
  hierarchy is the main point
\item
  the user needs focused keyboard inspection rather than passive comparison
\item
  the item set is effectively one-dimensional and would read more honestly as a list
\end{itemize}

Variants:

\begin{itemize}
\tightlist
\item
  \texttt{v\allowbreak{}a\allowbreak{}r\allowbreak{}i\allowbreak{}a\allowbreak{}n\allowbreak{}t\allowbreak{}:\allowbreak{} \allowbreak{}"\allowbreak{}b\allowbreak{}o\allowbreak{}x\allowbreak{}"} keeps the default Unicode boxed table
\item
  \texttt{v\allowbreak{}a\allowbreak{}r\allowbreak{}i\allowbreak{}a\allowbreak{}n\allowbreak{}t\allowbreak{}:\allowbreak{} \allowbreak{}"\allowbreak{}a\allowbreak{}s\allowbreak{}c\allowbreak{}i\allowbreak{}i\allowbreak{}-\allowbreak{}g\allowbreak{}r\allowbreak{}i\allowbreak{}d\allowbreak{}"} renders a portable \texttt{+\allowbreak{}-\allowbreak{}-\allowbreak{}-\allowbreak{}+} grid
\item
  \texttt{v\allowbreak{}a\allowbreak{}r\allowbreak{}i\allowbreak{}a\allowbreak{}n\allowbreak{}t\allowbreak{}:\allowbreak{} \allowbreak{}"\allowbreak{}r\allowbreak{}u\allowbreak{}l\allowbreak{}e\allowbreak{}d\allowbreak{}"} renders a borderless table with a heavy header rule and light row separators
\item
  \texttt{v\allowbreak{}a\allowbreak{}r\allowbreak{}i\allowbreak{}a\allowbreak{}n\allowbreak{}t\allowbreak{}:\allowbreak{} \allowbreak{}"\allowbreak{}h\allowbreak{}e\allowbreak{}a\allowbreak{}d\allowbreak{}e\allowbreak{}r\allowbreak{}-\allowbreak{}r\allowbreak{}u\allowbreak{}l\allowbreak{}e\allowbreak{}"} keeps only the header separator
\item
  \texttt{v\allowbreak{}a\allowbreak{}r\allowbreak{}i\allowbreak{}a\allowbreak{}n\allowbreak{}t\allowbreak{}:\allowbreak{} \allowbreak{}"\allowbreak{}p\allowbreak{}l\allowbreak{}a\allowbreak{}i\allowbreak{}n\allowbreak{}"} renders aligned columns with spaces and no rules
\item
  \texttt{v\allowbreak{}a\allowbreak{}r\allowbreak{}i\allowbreak{}a\allowbreak{}n\allowbreak{}t\allowbreak{}:\allowbreak{} \allowbreak{}"\allowbreak{}m\allowbreak{}a\allowbreak{}r\allowbreak{}k\allowbreak{}d\allowbreak{}o\allowbreak{}w\allowbreak{}n\allowbreak{}"} renders a GitHub-flavored Markdown pipe table
\item
  \texttt{v\allowbreak{}a\allowbreak{}r\allowbreak{}i\allowbreak{}a\allowbreak{}n\allowbreak{}t\allowbreak{}:\allowbreak{} \allowbreak{}"\allowbreak{}d\allowbreak{}e\allowbreak{}f\allowbreak{}i\allowbreak{}n\allowbreak{}i\allowbreak{}t\allowbreak{}i\allowbreak{}o\allowbreak{}n\allowbreak{}"} renders two-column field/value reports
\item
  \texttt{v\allowbreak{}a\allowbreak{}r\allowbreak{}i\allowbreak{}a\allowbreak{}n\allowbreak{}t\allowbreak{}:\allowbreak{} \allowbreak{}"\allowbreak{}e\allowbreak{}x\allowbreak{}p\allowbreak{}a\allowbreak{}n\allowbreak{}d\allowbreak{}e\allowbreak{}d\allowbreak{}"} renders psql-style record inspection
\end{itemize}

Examples:

\begin{CodeBlock}
box
┌───────┬────────┐
│ Name  │ State  │
├───────┼────────┤
│ Alice │ active │
└───────┴────────┘

ruled
Name   State
━━━━━  ━━━━━━
Alice  active

header-rule
Name   State
-----  ------
Alice  active

plain
Name   State
Alice  active

markdown
| Name  | State  |
|-------|--------|
| Alice | active |

definition
Field  Value
━━━━━  ━━━━━━
State  active

expanded
-[ RECORD 1 ]------------
Name  | Alice
State | active
\end{CodeBlock}

Lowering:

\begin{itemize}
\tightlist
\item
  \texttt{p\allowbreak{}i\allowbreak{}p\allowbreak{}e} mode still defaults to TSV
\item
  callers can opt into \texttt{p\allowbreak{}i\allowbreak{}p\allowbreak{}e\allowbreak{}F\allowbreak{}o\allowbreak{}r\allowbreak{}m\allowbreak{}a\allowbreak{}t\allowbreak{}:\allowbreak{} \allowbreak{}"\allowbreak{}c\allowbreak{}s\allowbreak{}v\allowbreak{}"}, \texttt{"\allowbreak{}m\allowbreak{}a\allowbreak{}r\allowbreak{}k\allowbreak{}d\allowbreak{}o\allowbreak{}w\allowbreak{}n\allowbreak{}"}, or \texttt{"\allowbreak{}a\allowbreak{}s\allowbreak{}c\allowbreak{}i\allowbreak{}i\allowbreak{}-\allowbreak{}g\allowbreak{}r\allowbreak{}i\allowbreak{}d\allowbreak{}"}
\item
  visual aligned-column variants use spaces, not tabs; tabs are reserved for TSV serialization
\end{itemize}

\subsection{\texorpdfstring{\texttt{t\allowbreak{}a\allowbreak{}b\allowbreak{}l\allowbreak{}e\allowbreak{}S\allowbreak{}u\allowbreak{}r\allowbreak{}f\allowbreak{}a\allowbreak{}c\allowbreak{}e\allowbreak{}(\allowbreak{})}}{}}\label{record-023-tablesurface}

Use when:

\begin{itemize}
\tightlist
\item
  the user needs the same passive comparison semantics as \texttt{t\allowbreak{}a\allowbreak{}b\allowbreak{}l\allowbreak{}e\allowbreak{}(\allowbreak{})}
\item
  the surrounding V3 runtime is already composing \texttt{S\allowbreak{}u\allowbreak{}r\allowbreak{}f\allowbreak{}a\allowbreak{}c\allowbreak{}e} output
\item
  cells may contain surface-native labels or badges
\end{itemize}

Avoid when:

\begin{itemize}
\tightlist
\item
  direct string output is the final destination
\item
  the app needs keyboard-owned row/cell focus semantics
\end{itemize}

\subsection{\texorpdfstring{\texttt{n\allowbreak{}a\allowbreak{}v\allowbreak{}i\allowbreak{}g\allowbreak{}a\allowbreak{}b\allowbreak{}l\allowbreak{}e\allowbreak{}T\allowbreak{}a\allowbreak{}b\allowbreak{}l\allowbreak{}e\allowbreak{}(\allowbreak{})} / \texttt{n\allowbreak{}a\allowbreak{}v\allowbreak{}i\allowbreak{}g\allowbreak{}a\allowbreak{}b\allowbreak{}l\allowbreak{}e\allowbreak{}T\allowbreak{}a\allowbreak{}b\allowbreak{}l\allowbreak{}e\allowbreak{}S\allowbreak{}u\allowbreak{}r\allowbreak{}f\allowbreak{}a\allowbreak{}c\allowbreak{}e\allowbreak{}(\allowbreak{})}}{ / }}\label{record-023-navigabletable-navigabletablesurface}

Use when:

\begin{itemize}
\tightlist
\item
  table data must be actively traversed
\item
  the app wants row/cell focus semantics
\item
  comparison across stable columns is still the job even while navigating
\end{itemize}

Avoid when:

\begin{itemize}
\tightlist
\item
  passive comparison is enough
\item
  the output must remain portable through core-only string rendering
\item
  the user is really selecting among one-dimensional records instead of comparing attributes
\end{itemize}

\subsection{\texorpdfstring{\texttt{b\allowbreak{}r\allowbreak{}o\allowbreak{}w\allowbreak{}s\allowbreak{}a\allowbreak{}b\allowbreak{}l\allowbreak{}e\allowbreak{}L\allowbreak{}i\allowbreak{}s\allowbreak{}t\allowbreak{}(\allowbreak{})} / \texttt{b\allowbreak{}r\allowbreak{}o\allowbreak{}w\allowbreak{}s\allowbreak{}a\allowbreak{}b\allowbreak{}l\allowbreak{}e\allowbreak{}L\allowbreak{}i\allowbreak{}s\allowbreak{}t\allowbreak{}S\allowbreak{}u\allowbreak{}r\allowbreak{}f\allowbreak{}a\allowbreak{}c\allowbreak{}e\allowbreak{}(\allowbreak{})}}{ / }}\label{record-023-browsablelist-browsablelistsurface}

Use when:

\begin{itemize}
\tightlist
\item
  items are one-dimensional
\item
  descriptions matter
\item
  the user is exploring or choosing among records
\end{itemize}

Avoid when:

\begin{itemize}
\tightlist
\item
  actions/commands are the result rather than records
\item
  column comparison is the point
\item
  hierarchy/path structure is the point
\end{itemize}

\subsection{\texorpdfstring{\texttt{f\allowbreak{}i\allowbreak{}l\allowbreak{}e\allowbreak{}P\allowbreak{}i\allowbreak{}c\allowbreak{}k\allowbreak{}e\allowbreak{}r\allowbreak{}(\allowbreak{})} / \texttt{f\allowbreak{}i\allowbreak{}l\allowbreak{}e\allowbreak{}P\allowbreak{}i\allowbreak{}c\allowbreak{}k\allowbreak{}e\allowbreak{}r\allowbreak{}S\allowbreak{}u\allowbreak{}r\allowbreak{}f\allowbreak{}a\allowbreak{}c\allowbreak{}e\allowbreak{}(\allowbreak{})}}{ / }}\label{record-023-filepicker-filepickersurface}

Use when:

\begin{itemize}
\tightlist
\item
  the user is traversing filesystem hierarchy
\item
  the current path and directory transitions are part of the job
\end{itemize}

Avoid when:

\begin{itemize}
\tightlist
\item
  the content is generic hierarchy rather than actual file/path navigation
\item
  actions/commands are the result rather than a selected file or directory
\end{itemize}

\section{3.5 Viewport masking and scroll ownership}\label{record-023-viewport-masking-and-scroll-ownership}

\subsection{Governing rules}\label{record-023-governing-rules}

\begin{itemize}
\tightlist
\item
  \texttt{v\allowbreak{}i\allowbreak{}e\allowbreak{}w\allowbreak{}p\allowbreak{}o\allowbreak{}r\allowbreak{}t\allowbreak{}S\allowbreak{}u\allowbreak{}r\allowbreak{}f\allowbreak{}a\allowbreak{}c\allowbreak{}e\allowbreak{}(\allowbreak{})} is the canonical masking primitive for bounded overflow scrolling in rich TUI composition.
\item
  \texttt{v\allowbreak{}i\allowbreak{}e\allowbreak{}w\allowbreak{}p\allowbreak{}o\allowbreak{}r\allowbreak{}t\allowbreak{}(\allowbreak{})} is the explicit text-lowering path when the content is intentionally string-first.
\item
  \texttt{p\allowbreak{}a\allowbreak{}g\allowbreak{}e\allowbreak{}r\allowbreak{}S\allowbreak{}u\allowbreak{}r\allowbreak{}f\allowbreak{}a\allowbreak{}c\allowbreak{}e\allowbreak{}(\allowbreak{})} and \texttt{f\allowbreak{}o\allowbreak{}c\allowbreak{}u\allowbreak{}s\allowbreak{}A\allowbreak{}r\allowbreak{}e\allowbreak{}a\allowbreak{}S\allowbreak{}u\allowbreak{}r\allowbreak{}f\allowbreak{}a\allowbreak{}c\allowbreak{}e\allowbreak{}(\allowbreak{})} are specialization layers on top of the viewport primitive, not separate scroll systems.
\item
  If a component only needs a clipped window plus scroll state, it should wrap viewport semantics rather than invent a bespoke \texttt{s\allowbreak{}c\allowbreak{}r\allowbreak{}o\allowbreak{}l\allowbreak{}l\allowbreak{}Y} slicer.
\item
  \texttt{n\allowbreak{}a\allowbreak{}v\allowbreak{}i\allowbreak{}g\allowbreak{}a\allowbreak{}b\allowbreak{}l\allowbreak{}e\allowbreak{}T\allowbreak{}a\allowbreak{}b\allowbreak{}l\allowbreak{}e\allowbreak{}S\allowbreak{}u\allowbreak{}r\allowbreak{}f\allowbreak{}a\allowbreak{}c\allowbreak{}e\allowbreak{}(\allowbreak{})} is the important exception: wrapped comparison rows keep row-aware scroll semantics, so generic line clipping is the wrong abstraction there.
\item
  Scroll ownership must stay unified. If keyboard focus, mouse hit-testing, and rendering disagree about the visible window height, the component has a layout bug.
\item
  Scrollbars belong to the viewport layer. Components should not fake overflow by moving rows around without giving the user a truthful masked window.
\end{itemize}

\subsection{Layout rules}\label{record-023-layout-rules}

\begin{itemize}
\tightlist
\item
  Components should either fill the rectangle their parent assigns them or be intentionally placed within it by a parent-owned layout helper such as \texttt{p\allowbreak{}l\allowbreak{}a\allowbreak{}c\allowbreak{}e\allowbreak{}S\allowbreak{}u\allowbreak{}r\allowbreak{}f\allowbreak{}a\allowbreak{}c\allowbreak{}e\allowbreak{}(\allowbreak{})}.
\item
  Parents own outer whitespace and anchoring. Child components should not silently shrink and leave accidental slack against a pane edge.
\item
  When a pane has a fixed header or separator row, the remaining rows should belong to one honest body region rather than an ad hoc mix of body plus invisible slack.
\item
  Overflow handling must be explicit: wrap, stack, clip, or viewport. Scrolling content without a matching viewport contract is not an acceptable fallback.
\end{itemize}

\subsection{\texorpdfstring{\texttt{v\allowbreak{}i\allowbreak{}e\allowbreak{}w\allowbreak{}p\allowbreak{}o\allowbreak{}r\allowbreak{}t\allowbreak{}S\allowbreak{}u\allowbreak{}r\allowbreak{}f\allowbreak{}a\allowbreak{}c\allowbreak{}e\allowbreak{}(\allowbreak{})}}{}}\label{record-023-viewportsurface}

Use when:

\begin{itemize}
\tightlist
\item
  an existing \texttt{S\allowbreak{}u\allowbreak{}r\allowbreak{}f\allowbreak{}a\allowbreak{}c\allowbreak{}e} or layout-backed region just needs a scrollable mask
\item
  the app already owns the surrounding chrome and only needs honest overflow scrolling
\item
  the component should stay on the structured \texttt{S\allowbreak{}u\allowbreak{}r\allowbreak{}f\allowbreak{}a\allowbreak{}c\allowbreak{}e} path
\end{itemize}

Avoid when:

\begin{itemize}
\tightlist
\item
  the content is intentionally string-first and the text-lowering path is the actual destination
\item
  the component needs domain-specific focus chrome or status semantics that belong in a higher-level wrapper
\end{itemize}

\section{3.6 Containment and formatted documents}\label{record-023-containment-and-formatted-documents}

\subsection{\texorpdfstring{\texttt{b\allowbreak{}o\allowbreak{}x\allowbreak{}(\allowbreak{})} / \texttt{h\allowbreak{}e\allowbreak{}a\allowbreak{}d\allowbreak{}e\allowbreak{}r\allowbreak{}B\allowbreak{}o\allowbreak{}x\allowbreak{}(\allowbreak{})}}{ / }}\label{record-023-box-headerbox}

Use when:

\begin{itemize}
\tightlist
\item
  a region needs visible containment or a local title
\item
  peer panels need to feel like separate working areas
\item
  compact metadata belongs in a header rather than in the surrounding prose
\end{itemize}

Avoid when:

\begin{itemize}
\tightlist
\item
  every subsection would get a border just for visual decoration
\item
  the real job is status escalation, not containment
\item
  a divider, heading, or simple whitespace would communicate structure more honestly
\end{itemize}

\subsection{\texorpdfstring{\texttt{m\allowbreak{}a\allowbreak{}r\allowbreak{}k\allowbreak{}d\allowbreak{}o\allowbreak{}w\allowbreak{}n\allowbreak{}(\allowbreak{})}}{}}\label{record-023-markdown}

Use when:

\begin{itemize}
\tightlist
\item
  the content is help, reference, release notes, or bounded prose
\item
  lightweight structure such as headings, lists, quotes, links, and code fences materially helps comprehension
\item
  the same content should lower honestly across rich, pipe, and accessible modes
\end{itemize}

Avoid when:

\begin{itemize}
\tightlist
\item
  the content is so large that it really needs document navigation instead of one rendered block
\item
  the prose is user-authored rich content that expects full browser-grade markdown semantics
\item
  the app is really laying out UI regions or controls and is trying to use markdown as a layout engine
\item
  dense comparison, forms, or command discovery are the real jobs
\end{itemize}

\subsection{\texorpdfstring{\texttt{h\allowbreak{}y\allowbreak{}p\allowbreak{}e\allowbreak{}r\allowbreak{}l\allowbreak{}i\allowbreak{}n\allowbreak{}k\allowbreak{}(\allowbreak{})}}{}}\label{record-023-hyperlink}

Use when:

\begin{itemize}
\tightlist
\item
  the destination itself should remain part of the rendered output
\item
  terminals with link support can benefit from clickability without hiding where the link goes
\end{itemize}

Avoid when:

\begin{itemize}
\tightlist
\item
  the label is vague or generic
\item
  trust depends on seeing the actual destination, but the surrounding content hides it
\item
  the interaction is really an app-owned action instead of an external destination
\end{itemize}

\subsection{\texorpdfstring{\texttt{k\allowbreak{}b\allowbreak{}d\allowbreak{}(\allowbreak{})}}{}}\label{record-023-kbd}

Use when:

\begin{itemize}
\tightlist
\item
  a local action needs an inline shortcut cue
\item
  the key hint belongs right next to the control, field, or instruction it affects
\end{itemize}

Avoid when:

\begin{itemize}
\tightlist
\item
  the user needs a grouped command reference or shell-level keymap
\item
  shortcut chips would start competing with the real content for attention
\end{itemize}

\subsection{\texorpdfstring{\texttt{p\allowbreak{}a\allowbreak{}g\allowbreak{}e\allowbreak{}r\allowbreak{}S\allowbreak{}u\allowbreak{}r\allowbreak{}f\allowbreak{}a\allowbreak{}c\allowbreak{}e\allowbreak{}(\allowbreak{})} / \texttt{f\allowbreak{}o\allowbreak{}c\allowbreak{}u\allowbreak{}s\allowbreak{}A\allowbreak{}r\allowbreak{}e\allowbreak{}a\allowbreak{}S\allowbreak{}u\allowbreak{}r\allowbreak{}f\allowbreak{}a\allowbreak{}c\allowbreak{}e\allowbreak{}(\allowbreak{})}}{ / }}\label{record-023-pagersurface-focusareasurface}

Use when:

\begin{itemize}
\tightlist
\item
  the scrollable region needs explicit line status, current-position framing, or focus gutter semantics
\item
  the pane is part of a larger workspace that needs visible ownership/focus signals
\end{itemize}

Avoid when:

\begin{itemize}
\tightlist
\item
  a plain masked region is enough and extra chrome would distract from the task
\end{itemize}

\subsection{Choosing among browseable structures}\label{record-023-choosing-among-browseable-structures}

\begin{itemize}
\tightlist
\item
  Choose \texttt{b\allowbreak{}r\allowbreak{}o\allowbreak{}w\allowbreak{}s\allowbreak{}a\allowbreak{}b\allowbreak{}l\allowbreak{}e\allowbreak{}L\allowbreak{}i\allowbreak{}s\allowbreak{}t\allowbreak{}S\allowbreak{}u\allowbreak{}r\allowbreak{}f\allowbreak{}a\allowbreak{}c\allowbreak{}e\allowbreak{}(\allowbreak{})} when rows are essentially labels plus brief descriptions.
\item
  Choose \texttt{c\allowbreak{}o\allowbreak{}m\allowbreak{}m\allowbreak{}a\allowbreak{}n\allowbreak{}d\allowbreak{}P\allowbreak{}a\allowbreak{}l\allowbreak{}e\allowbreak{}t\allowbreak{}t\allowbreak{}e\allowbreak{}S\allowbreak{}u\allowbreak{}r\allowbreak{}f\allowbreak{}a\allowbreak{}c\allowbreak{}e\allowbreak{}(\allowbreak{})} when the outcome is an action or navigation command.
\item
  Choose \texttt{f\allowbreak{}i\allowbreak{}l\allowbreak{}e\allowbreak{}P\allowbreak{}i\allowbreak{}c\allowbreak{}k\allowbreak{}e\allowbreak{}r\allowbreak{}S\allowbreak{}u\allowbreak{}r\allowbreak{}f\allowbreak{}a\allowbreak{}c\allowbreak{}e\allowbreak{}(\allowbreak{})} when path traversal is the mental model.
\item
  Choose \texttt{n\allowbreak{}a\allowbreak{}v\allowbreak{}i\allowbreak{}g\allowbreak{}a\allowbreak{}b\allowbreak{}l\allowbreak{}e\allowbreak{}T\allowbreak{}a\allowbreak{}b\allowbreak{}l\allowbreak{}e\allowbreak{}S\allowbreak{}u\allowbreak{}r\allowbreak{}f\allowbreak{}a\allowbreak{}c\allowbreak{}e\allowbreak{}(\allowbreak{})} when comparison across columns remains the job even while navigating.
\item
  Choose \texttt{v\allowbreak{}i\allowbreak{}e\allowbreak{}w\allowbreak{}p\allowbreak{}o\allowbreak{}r\allowbreak{}t\allowbreak{}S\allowbreak{}u\allowbreak{}r\allowbreak{}f\allowbreak{}a\allowbreak{}c\allowbreak{}e\allowbreak{}(\allowbreak{})} directly when the content already exists as a region and just needs bounded overflow masking.
\end{itemize}

\subsection{\texorpdfstring{\texttt{t\allowbreak{}r\allowbreak{}e\allowbreak{}e\allowbreak{}(\allowbreak{})}}{}}\label{record-023-tree}

Use when:

\begin{itemize}
\tightlist
\item
  the primary relationship is parent/child
\item
  the user needs to understand nesting or containment
\end{itemize}

Avoid when:

\begin{itemize}
\tightlist
\item
  multiple parents or dependency edges matter
\item
  the user is really traversing filesystem paths, where \texttt{f\allowbreak{}i\allowbreak{}l\allowbreak{}e\allowbreak{}P\allowbreak{}i\allowbreak{}c\allowbreak{}k\allowbreak{}e\allowbreak{}r\allowbreak{}S\allowbreak{}u\allowbreak{}r\allowbreak{}f\allowbreak{}a\allowbreak{}c\allowbreak{}e\allowbreak{}(\allowbreak{})} is more honest
\item
  dense row/column comparison would communicate the content more clearly than nesting
\end{itemize}

\subsection{\texorpdfstring{\texttt{d\allowbreak{}a\allowbreak{}g\allowbreak{}(\allowbreak{})} / \texttt{d\allowbreak{}a\allowbreak{}g\allowbreak{}P\allowbreak{}a\allowbreak{}n\allowbreak{}e\allowbreak{}(\allowbreak{})}}{ / }}\label{record-023-dag-dagpane}

Use when:

\begin{itemize}
\tightlist
\item
  the primary relationship is dependency or flow
\item
  multiple parents or graph edges matter
\item
  the user needs to answer questions like ``what does this depend on?'' or ``what does this unblock?''
\end{itemize}

Avoid when:

\begin{itemize}
\tightlist
\item
  a tree, timeline, or table would answer the question more directly
\item
  the graph would mostly act as architecture wallpaper instead of supporting a concrete task
\item
  the user really needs a metrics summary or focused fragment instead of the whole graph
\end{itemize}

\subsection{\texorpdfstring{\texttt{t\allowbreak{}i\allowbreak{}m\allowbreak{}e\allowbreak{}l\allowbreak{}i\allowbreak{}n\allowbreak{}e\allowbreak{}(\allowbreak{})}}{}}\label{record-023-timeline}

Use when:

\begin{itemize}
\tightlist
\item
  chronology is the main structure
\item
  events should read as a sequence, milestone trail, or audit-oriented progression
\end{itemize}

Avoid when:

\begin{itemize}
\tightlist
\item
  dependency or causality matters more than ordering
\item
  the content is so dense that a table, log, or summary-plus-detail drill-down would be more readable
\item
  the user is really comparing attributes across many events rather than following a temporal story
\end{itemize}

\section{3.7 Expressive branding and decorative emphasis}\label{record-023-expressive-branding-and-decorative-emphasis}

\subsection{\texorpdfstring{\texttt{l\allowbreak{}o\allowbreak{}a\allowbreak{}d\allowbreak{}R\allowbreak{}a\allowbreak{}n\allowbreak{}d\allowbreak{}o\allowbreak{}m\allowbreak{}L\allowbreak{}o\allowbreak{}g\allowbreak{}o\allowbreak{}(\allowbreak{})} / \texttt{g\allowbreak{}r\allowbreak{}a\allowbreak{}d\allowbreak{}i\allowbreak{}e\allowbreak{}n\allowbreak{}t\allowbreak{}T\allowbreak{}e\allowbreak{}x\allowbreak{}t\allowbreak{}(\allowbreak{})}}{ / }}\label{record-023-loadrandomlogo-gradienttext}

Use when:

\begin{itemize}
\tightlist
\item
  the product needs a deliberate branded or celebratory moment
\item
  splash surfaces, docs, or demos benefit from a little atmosphere
\end{itemize}

Avoid when:

\begin{itemize}
\tightlist
\item
  ordinary workspace chrome, navigation, status, or instructions need to stay neutral and legible
\item
  the emphasis would compete with the actual task
\item
  the color treatment is being asked to carry meaning that should live in text
\end{itemize}

\section{5. App-authored primitives}\label{record-023-app-authored-primitives}

\subsection{\texorpdfstring{\texttt{r\allowbreak{}e\allowbreak{}n\allowbreak{}d\allowbreak{}e\allowbreak{}r\allowbreak{}B\allowbreak{}y\allowbreak{}M\allowbreak{}o\allowbreak{}d\allowbreak{}e\allowbreak{}(\allowbreak{})}}{}}\label{record-023-renderbymode}

Use when:

\begin{itemize}
\tightlist
\item
  the app needs a domain-specific primitive that the shared component library should not own
\item
  the same semantic concept must lower honestly across rich, pipe, and accessible output
\end{itemize}

Avoid when:

\begin{itemize}
\tightlist
\item
  a shipped Bijou family already matches the job
\item
  the branching only exists to chase visual novelty instead of preserving meaning
\end{itemize}

\section{4. Disclosure and progressive complexity}\label{record-023-disclosure-and-progressive-complexity}

\subsection{\texorpdfstring{\texttt{a\allowbreak{}c\allowbreak{}c\allowbreak{}o\allowbreak{}r\allowbreak{}d\allowbreak{}i\allowbreak{}o\allowbreak{}n\allowbreak{}(\allowbreak{})}}{}}\label{record-023-accordion}

Use when:

\begin{itemize}
\tightlist
\item
  content is secondary or optional
\item
  one page can hold both summary and detail
\end{itemize}

Avoid when:

\begin{itemize}
\tightlist
\item
  the sections are peer destinations that deserve explicit navigation
\end{itemize}

\subsection{\texorpdfstring{\texttt{t\allowbreak{}a\allowbreak{}b\allowbreak{}s\allowbreak{}(\allowbreak{})}}{}}\label{record-023-tabs}

Use when:

\begin{itemize}
\tightlist
\item
  views are peers
\item
  only one view should be foregrounded at a time
\end{itemize}

Avoid when:

\begin{itemize}
\tightlist
\item
  the user is progressing through steps
\item
  hierarchy/path context matters more than peer switching
\end{itemize}

\subsection{\texorpdfstring{\texttt{s\allowbreak{}t\allowbreak{}e\allowbreak{}p\allowbreak{}p\allowbreak{}e\allowbreak{}r\allowbreak{}(\allowbreak{})}}{}}\label{record-023-stepper}

Use when:

\begin{itemize}
\tightlist
\item
  the user is moving through a process
\item
  order matters
\item
  the current stage should be understandable at a glance
\end{itemize}

Avoid when:

\begin{itemize}
\tightlist
\item
  the views are non-sequential
\item
  the user is really switching among peer destinations instead of progressing through a flow
\end{itemize}

\subsection{\texorpdfstring{\texttt{b\allowbreak{}r\allowbreak{}e\allowbreak{}a\allowbreak{}d\allowbreak{}c\allowbreak{}r\allowbreak{}u\allowbreak{}m\allowbreak{}b\allowbreak{}(\allowbreak{})}}{}}\label{record-023-breadcrumb}

Use when:

\begin{itemize}
\tightlist
\item
  the user needs path context
\item
  hierarchy or nesting should stay visible while reading the current location
\end{itemize}

Avoid when:

\begin{itemize}
\tightlist
\item
  peer destinations are the real model
\item
  the path is so deep or unstable that a simpler label plus local navigation would be clearer
\end{itemize}

\subsection{\texorpdfstring{\texttt{p\allowbreak{}a\allowbreak{}g\allowbreak{}i\allowbreak{}n\allowbreak{}a\allowbreak{}t\allowbreak{}o\allowbreak{}r\allowbreak{}(\allowbreak{})}}{}}\label{record-023-paginator}

Use when:

\begin{itemize}
\tightlist
\item
  the current page or position within a bounded sequence matters
\item
  page count should be visible without turning into a full navigation menu
\end{itemize}

Avoid when:

\begin{itemize}
\tightlist
\item
  the user needs rich peer navigation rather than compact progress
\item
  the content is an infinite feed or unbounded stream where page counts are misleading
\end{itemize}

\section{5. Overlay escalation}\label{record-023-overlay-escalation}

Choose overlays by how much they should interrupt and how much context they should retain.

\subsection{Governing rules}\label{record-023-governing-rules-1}

\begin{itemize}
\tightlist
\item
  Overlays are for interruption, review, or supplemental context. They are not substitutes for normal page structure.
\item
  Pick the least interruptive overlay that still preserves the meaning and action path.
\item
  If users need history, stacking, routing, or recall, move up from \texttt{t\allowbreak{}o\allowbreak{}a\allowbreak{}s\allowbreak{}t\allowbreak{}(\allowbreak{})} to the notification system.
\item
  If background input must stop, use \texttt{m\allowbreak{}o\allowbreak{}d\allowbreak{}a\allowbreak{}l\allowbreak{}(\allowbreak{})} rather than trying to make a toast or drawer behave like one.
\item
  If the content is long-lived and belongs in the reading flow, use \texttt{a\allowbreak{}l\allowbreak{}e\allowbreak{}r\allowbreak{}t\allowbreak{}(\allowbreak{})} or a normal page region instead of any overlay.
\item
  Use structured \texttt{S\allowbreak{}u\allowbreak{}r\allowbreak{}f\allowbreak{}a\allowbreak{}c\allowbreak{}e} content inside overlays when the interrupting layer needs real rows, embedded components, or richer layout. Flatten to plain text only at explicit lowering boundaries.
\end{itemize}

\subsection{\texorpdfstring{\texttt{t\allowbreak{}o\allowbreak{}o\allowbreak{}l\allowbreak{}t\allowbreak{}i\allowbreak{}p\allowbreak{}(\allowbreak{})}}{}}\label{record-023-tooltip}

\begin{itemize}
\tightlist
\item
  for tiny local explanation attached to nearby content
\item
  prefer one or two short lines that clarify a control, label, or state
\item
  not for decisions
\item
  not for content that needs scrolling, recall, commands, or workflow ownership
\end{itemize}

\subsection{\texorpdfstring{\texttt{d\allowbreak{}r\allowbreak{}a\allowbreak{}w\allowbreak{}e\allowbreak{}r\allowbreak{}(\allowbreak{})}}{}}\label{record-023-drawer}

\begin{itemize}
\tightlist
\item
  for supplemental context, inspectors, side work
\item
  keep the main surface visible
\item
  prefer when the user needs reference material while still working in the main task
\item
  use panel-scoped drawers when the supplemental context belongs to one workspace region instead of the whole frame
\item
  do not turn a drawer into a modal by stuffing it with blocking confirmation copy
\end{itemize}

\subsection{\texorpdfstring{\texttt{m\allowbreak{}o\allowbreak{}d\allowbreak{}a\allowbreak{}l\allowbreak{}(\allowbreak{})}}{}}\label{record-023-modal-1}

\begin{itemize}
\tightlist
\item
  for blocking review, choice, or confirmation
\item
  background shortcuts and pointer actions should be blocked
\item
  use sparingly; interruption must be justified
\item
  prefer short review-and-decide tasks over deep browsing or long-lived side work
\end{itemize}

\subsection{\texorpdfstring{\texttt{t\allowbreak{}o\allowbreak{}a\allowbreak{}s\allowbreak{}t\allowbreak{}(\allowbreak{})}}{}}\label{record-023-toast-1}

\begin{itemize}
\tightlist
\item
  for one transient overlay the app is composing directly
\item
  not for stack/history/routing problems
\item
  if the same class of event happens repeatedly, stop composing ad hoc toasts and move up to notifications
\end{itemize}

\subsection{notifications}\label{record-023-notifications-1}

\begin{itemize}
\tightlist
\item
  for events, not documents
\item
  use history/archive when recall matters
\item
  prefer over repeated ad hoc toasts once app-level messaging becomes systemic
\item
  route app-owned warnings and errors here when users may need to revisit them after the initial interruption
\item
  keep actions singular and obvious; do not hide multi-step workflows inside notification cards
\end{itemize}

\section{6. Forms and progression}\label{record-023-forms-and-progression}

\subsection{\texorpdfstring{\texttt{g\allowbreak{}r\allowbreak{}o\allowbreak{}u\allowbreak{}p\allowbreak{}(\allowbreak{})}}{}}\label{record-023-group}

Use when:

\begin{itemize}
\tightlist
\item
  the form fits one mental step
\item
  fields are directly related
\end{itemize}

\subsection{\texorpdfstring{\texttt{w\allowbreak{}i\allowbreak{}z\allowbreak{}a\allowbreak{}r\allowbreak{}d\allowbreak{}(\allowbreak{})}}{}}\label{record-023-wizard}

Use when:

\begin{itemize}
\tightlist
\item
  the flow has stages
\item
  branching or transformation matters
\item
  the user benefits from progressive disclosure
\end{itemize}

\subsection{\texorpdfstring{\texttt{c\allowbreak{}o\allowbreak{}n\allowbreak{}f\allowbreak{}i\allowbreak{}r\allowbreak{}m\allowbreak{}(\allowbreak{})}}{}}\label{record-023-confirm}

Use when:

\begin{itemize}
\tightlist
\item
  the decision is truly binary
\end{itemize}

Avoid when:

\begin{itemize}
\tightlist
\item
  the user needs to compare options or consequences in detail
\end{itemize}

\section{7. App shell composition}\label{record-023-app-shell-composition}

\subsection{Governing rules}\label{record-023-governing-rules-2}

\begin{itemize}
\tightlist
\item
  The shell exists to frame destinations, workspace state, and global actions. It should not compete with page content for narrative space.
\item
  Use shell chrome for orientation, status, and cross-cutting actions. Put explanation and task detail back into the page.
\item
  Prefer the simplest workspace that preserves the mental model. Do not use split panes or grids just because the runtime can.
\item
  If a region is not meaningfully inspectable, comparable, or supplemental, it probably should not be its own pane.
\item
  Keep shell roles distinct:

  \begin{itemize}
  \tightlist
  \item
    status rails communicate concise global state
  \item
    help hints teach shortcuts and scope
  \item
    command palette surfaces action discovery
  \item
    notifications surface events and follow-up
  \end{itemize}
\item
  If one shell element starts doing another shell element's job, the design is drifting.
\end{itemize}

\subsection{\texorpdfstring{\texttt{c\allowbreak{}r\allowbreak{}e\allowbreak{}a\allowbreak{}t\allowbreak{}e\allowbreak{}F\allowbreak{}r\allowbreak{}a\allowbreak{}m\allowbreak{}e\allowbreak{}d\allowbreak{}A\allowbreak{}p\allowbreak{}p\allowbreak{}(\allowbreak{})}}{}}\label{record-023-createframedapp}

Use when:

\begin{itemize}
\tightlist
\item
  the app has multiple views or work areas
\item
  shell chrome, palette, overlays, and notifications should be standardized
\item
  the frame should own the default hosted runtime path via \texttt{a\allowbreak{}p\allowbreak{}p\allowbreak{}.\allowbreak{}r\allowbreak{}u\allowbreak{}n\allowbreak{}(\allowbreak{}.\allowbreak{}.\allowbreak{}.\allowbreak{})} or \texttt{r\allowbreak{}u\allowbreak{}n\allowbreak{}F\allowbreak{}r\allowbreak{}a\allowbreak{}m\allowbreak{}e\allowbreak{}d\allowbreak{}A\allowbreak{}p\allowbreak{}p\allowbreak{}(\allowbreak{}.\allowbreak{}.\allowbreak{}.\allowbreak{})}
\end{itemize}

Avoid when:

\begin{itemize}
\tightlist
\item
  the app is really just one screen or one prompt loop
\end{itemize}

Content guidance:

\begin{itemize}
\tightlist
\item
  tabs should represent peer destinations, not unrelated commands
\item
  shell notifications and overlays should remain subordinate to the active page task
\item
  help text should clarify shell behavior, not restate page content
\item
  command palette entries should prefer actions and navigation targets, not field-style data entry or explanatory prose
\item
  use drawers for supplemental workspace context and notifications for event messaging; do not overload the status rail with either
\item
  prefer the self-running framed-shell path before dropping down to raw \texttt{r\allowbreak{}u\allowbreak{}n\allowbreak{}(\allowbreak{}a\allowbreak{}p\allowbreak{}p\allowbreak{})}
\end{itemize}

\subsection{\texorpdfstring{\texttt{s\allowbreak{}t\allowbreak{}a\allowbreak{}t\allowbreak{}u\allowbreak{}s\allowbreak{}B\allowbreak{}a\allowbreak{}r\allowbreak{}(\allowbreak{})} and \texttt{s\allowbreak{}t\allowbreak{}a\allowbreak{}t\allowbreak{}u\allowbreak{}s\allowbreak{}B\allowbreak{}a\allowbreak{}r\allowbreak{}S\allowbreak{}u\allowbreak{}r\allowbreak{}f\allowbreak{}a\allowbreak{}c\allowbreak{}e\allowbreak{}(\allowbreak{})}}{ and }}\label{record-023-statusbar-and-statusbarsurface}

Use when:

\begin{itemize}
\tightlist
\item
  global state, mode, or shortcut context belongs at the shell edge
\end{itemize}

Avoid when:

\begin{itemize}
\tightlist
\item
  the content is explanatory or multi-line
\end{itemize}

Content guidance:

\begin{itemize}
\tightlist
\item
  keep status text short, global, and low-drama
\item
  use \texttt{s\allowbreak{}t\allowbreak{}a\allowbreak{}t\allowbreak{}u\allowbreak{}s\allowbreak{}B\allowbreak{}a\allowbreak{}r\allowbreak{}S\allowbreak{}u\allowbreak{}r\allowbreak{}f\allowbreak{}a\allowbreak{}c\allowbreak{}e\allowbreak{}(\allowbreak{})} when shell chrome is already being composed as a \texttt{S\allowbreak{}u\allowbreak{}r\allowbreak{}f\allowbreak{}a\allowbreak{}c\allowbreak{}e}
\item
  keep \texttt{s\allowbreak{}t\allowbreak{}a\allowbreak{}t\allowbreak{}u\allowbreak{}s\allowbreak{}B\allowbreak{}a\allowbreak{}r\allowbreak{}(\allowbreak{})} for explicit text output or lowering
\item
  reserve strong emphasis for mode changes, faults, or focus-critical state
\item
  do not turn the status bar into a secondary page body
\item
  do not route transient warnings or task-specific explanation here when notifications or in-page content would be more honest
\end{itemize}

\subsection{\texorpdfstring{\texttt{h\allowbreak{}e\allowbreak{}l\allowbreak{}p\allowbreak{}S\allowbreak{}h\allowbreak{}o\allowbreak{}r\allowbreak{}t\allowbreak{}(\allowbreak{})} / \texttt{h\allowbreak{}e\allowbreak{}l\allowbreak{}p\allowbreak{}V\allowbreak{}i\allowbreak{}e\allowbreak{}w\allowbreak{}(\allowbreak{})} and surface companions}{ /  and surface companions}}\label{record-023-helpshort-helpview-and-surface-companions}

Use when:

\begin{itemize}
\tightlist
\item
  keyboard ownership is real and shortcut discovery needs a concise hint or grouped reference
\end{itemize}

Avoid when:

\begin{itemize}
\tightlist
\item
  the text would only restate visible controls without clarifying behavior or scope
\end{itemize}

Content guidance:

\begin{itemize}
\tightlist
\item
  use \texttt{h\allowbreak{}e\allowbreak{}l\allowbreak{}p\allowbreak{}S\allowbreak{}h\allowbreak{}o\allowbreak{}r\allowbreak{}t\allowbreak{}S\allowbreak{}u\allowbreak{}r\allowbreak{}f\allowbreak{}a\allowbreak{}c\allowbreak{}e\allowbreak{}(\allowbreak{})} for shell hints that stay on the structured \texttt{S\allowbreak{}u\allowbreak{}r\allowbreak{}f\allowbreak{}a\allowbreak{}c\allowbreak{}e} path
\item
  use \texttt{h\allowbreak{}e\allowbreak{}l\allowbreak{}p\allowbreak{}V\allowbreak{}i\allowbreak{}e\allowbreak{}w\allowbreak{}S\allowbreak{}u\allowbreak{}r\allowbreak{}f\allowbreak{}a\allowbreak{}c\allowbreak{}e\allowbreak{}(\allowbreak{})} for grouped shortcut references embedded inside a rich TUI surface or modal
\item
  keep the plain string helpers for explicit lowering, pipe output, or textual docs/examples
\item
  group by user intent (\texttt{N\allowbreak{}a\allowbreak{}v\allowbreak{}i\allowbreak{}g\allowbreak{}a\allowbreak{}t\allowbreak{}i\allowbreak{}o\allowbreak{}n}, \texttt{A\allowbreak{}c\allowbreak{}t\allowbreak{}i\allowbreak{}o\allowbreak{}n\allowbreak{}s}, \texttt{S\allowbreak{}e\allowbreak{}l\allowbreak{}e\allowbreak{}c\allowbreak{}t\allowbreak{}i\allowbreak{}o\allowbreak{}n}) rather than by key shape
\item
  help is for shortcut discovery and scope clarification, not for executing commands or replacing a command palette
\end{itemize}

\subsection{\texorpdfstring{\texttt{s\allowbreak{}p\allowbreak{}l\allowbreak{}i\allowbreak{}t\allowbreak{}P\allowbreak{}a\allowbreak{}n\allowbreak{}e\allowbreak{}(\allowbreak{})} and \texttt{g\allowbreak{}r\allowbreak{}i\allowbreak{}d\allowbreak{}(\allowbreak{})}}{ and }}\label{record-023-splitpane-and-grid}

Use when:

\begin{itemize}
\tightlist
\item
  spatial composition helps the user reason about the workspace
\end{itemize}

Avoid when:

\begin{itemize}
\tightlist
\item
  the layout would be simpler and more legible as a sequential flow
\end{itemize}

Content guidance:

\begin{itemize}
\tightlist
\item
  use split panes for comparison, side-by-side reference, or inspector-style secondary context
\item
  use grids when the user is reading across stable regions, not when one long column would be clearer
\item
  label regions by job, not by visual position alone
\end{itemize}

\section{8. Mouse and multimodal interaction}\label{record-023-mouse-and-multimodal-interaction}

Mouse support in Bijou should follow these rules:

\begin{itemize}
\tightlist
\item
  keyboard is baseline
\item
  mouse is enhancement
\item
  every click target should mirror an existing keyboard path
\item
  overlays consume pointer input before background content
\item
  drag is reserved for spatial interactions such as pane resizing
\end{itemize}

Good early mouse targets:

\begin{itemize}
\tightlist
\item
  tabs
\item
  split dividers
\item
  notification actions and dismiss buttons
\item
  list/table row focus
\end{itemize}

Bad early mouse targets:

\begin{itemize}
\tightlist
\item
  hover-only critical meaning
\item
  interactions that have no keyboard parity
\end{itemize}

\section{9. Lowering across modes}\label{record-023-lowering-across-modes}

Every pattern should have a graceful lowering story.

\subsection{Examples}\label{record-023-examples}

\section{10. Motion and shader effects}\label{record-023-motion-and-shader-effects}

Motion and shader effects are supporting patterns, not the product itself.

Use them when:

\begin{itemize}
\tightlist
\item
  they clarify transition, progress, focus, or spatial change
\item
  the user benefits from seeing continuity rather than a sudden visual jump
\item
  the effect reinforces the information architecture instead of competing with it
\end{itemize}

Avoid them when:

\begin{itemize}
\tightlist
\item
  they are only decorative noise
\item
  they make dense work harder to read
\item
  the same information would be clearer as stable content, a status change, or a simpler transition
\end{itemize}

Specific guidance:

\begin{itemize}
\item
  use transition shaders to reinforce page or workspace change, not to show off that the renderer can animate
\item
  use \texttt{c\allowbreak{}a\allowbreak{}n\allowbreak{}v\allowbreak{}a\allowbreak{}s\allowbreak{}(\allowbreak{})} for purposeful visual surfaces such as splash, background atmosphere, or dedicated visual moments, not as a default wrapper around ordinary productivity screens
\item
  reduced-motion and non-interactive modes must still tell the truth even when the effect disappears entirely
\item
  notifications may lower to logs or in-flow alerts
\item
  command palette may lower to documented shortcuts or static command lists
\item
  navigable inspection views may lower to passive read-only views
\item
  modals may lower to step-by-step prompts or inline confirmation blocks
\end{itemize}

The docs for each family should make that lowering explicit instead of assuming an interactive full-screen terminal.

